
\documentclass[12pt]{article}

\addtolength{\textwidth}{.8in}
\addtolength{\hoffset}{-.4in}
\addtolength{\textheight}{1in}
\addtolength{\voffset}{-1in}

\usepackage{amsmath}
\usepackage{graphicx}
\usepackage{upquote} % to convert funny quotes to straight quotes
\usepackage{amsfonts}
\usepackage{multicol}
\usepackage{paralist}
\setlength{\parindent}{15pt}

\renewcommand*{\thefootnote}{\fnsymbol{footnote}}


\begin{document}

\begin{center}
\begin{Large}
DATA 219 Notes -- Apr. 5, 2018
\end{Large}
\end{center}

\section{k-Nearest Neighbors classifiers}

Another common approach to the ML classification problem is sort of at the
opposite extreme from Naive Bayes. Instead of very generally assuming that
each of the features is an independent and separate indicator of a data
point's class, the \textbf{kNN} (for ``$k$-nearest neighbors") algorithm
assumes that a point's features are very much bound together. In order to
classify a point, a kNN classifier finds the \textit{closest} training points
to it, and predicts the class that those points suggest (majority rules).

The parameter $k$ in kNN is simply the number of ``closest training points"
we'll consider. It is nearly always an odd number, so that in the two-class
case there are no ties. Suppose we have a 7-NN classifier, and we're asked to
classify a particular person as either \texttt{Jedi} or \texttt{Sith}. If we
find the seven most similar training points to our person, and discover that
four of them are \texttt{Jedi} and the other three are \texttt{Sith}, we'll
predict ``\texttt{Jedi}." The basic idea is that simple.

\subsection{Distance measures}

Much hinges on the question of \textit{closeness}. What do we mean when we say
two data points are ``close" together? There are some options here, many of
which are common sense. The first approach you might think of is to take the
standard \textit{Euclidean distance}, also called the ``straight-line
distance," between the data points. We just use the Pythagorean Theorem where
each feature is plotted on an axis. And in fact this is the most common
choice. Remember that the distance between two $n$-dimensional points $x$ and
$y$ is:

\begin{align*}
d = \sqrt{(x_1-y_1)^2 + (x_2-y_2)^2 + (x_3-y_3)^2 + \dots + (x_n-y_n)^2}
\end{align*}

In other words, the straightforward square-root-of-sum-of-squared-differences
works in \textit{any} number of dimensions, not just two.

This isn't the only approach people use. In general, we can use any
\textbf{distance measure} (or ``distance function," or ``distance metric")
that assigns a numerical ``distance" value to any two points. To be valid, a
distance measure must simply satisfy certain common sense
rules.\footnote{Including: the distance between two points must always be
greater than or equal to 0; the distance from $x$ to $y$ must be the same as
the distance from $y$ to $x$; and the distance from $x$ to $z$ must not be any
greater than the distance from $x$ to $y$ to $z$ (the ``triangle
inequality.")} Another common choice is the \textbf{Manhattan distance} (or
``taxicab distance") which is essentially the shortest path between the two
points that doesn't travel diagonally (\textit{i.e.}, the path only travels
along axes.)

\subsection{Feature standardization}

One very common issue is data sets that have features on widely different
scales. Suppose one of our features is a person's age, and another is that
person's salary. If we take the straightforward Euclidean distance between two
people, we'll find that the salary difference -- measured in the tens of
thousands -- completely overwhelms the age difference, which is at most a few
decades. Clearly a difference of \$50/year in salary isn't nearly as
significant as a difference of 50 years in age!

The solution is to \textbf{transform} our features so they are all
\textbf{standardized} (or ``Z-score normalized.") All we have to do is take
each feature, and:

\begin{compactenum}
\item subtract its mean from each point, then
\item divide each point by its standard deviation.
\end{compactenum}

This puts all features on the same level playing field, so to speak. If you
remember from your stats class, what we've done is compute the
\textbf{z-score} for each data point: the number of standard deviations above
(or below) the mean it is.

\subsection{Categorical features}

What do we do if our features are categorical? Same thing we learned in the
Naive Bayes lesson: make them numerical, in a principled way. Again, the most
common approach by far is to ``one-hot" encode each categorical feature as a
set of binary features, one for each feature value. (Recall that for
lightsabers, this would turn the single feature value ``\texttt{red}" into a
set of three features ``\texttt{0}," ``\texttt{0}," and ``\texttt{1}," meaning
``this data point has zero ``blueness," zero ``greenness," and one
``redness.")

You can use \texttt{np.where()} function to create one-hot vectors, or, use
the scikit-learn package's helpful ``\texttt{LabelBinarizer.fit\_transform}
function:

\small
\begin{Verbatim}[fontsize=\small,samepage=true,frame=none]
from sklearn.preprocessing import LabelBinarizer
lb = LabelBinarizer()
lightsaber = lb.fit_transform(df.lightsaber)
\end{Verbatim}
\normalsize

Careful: although this function takes a Series, it returns a NumPy array:

\small
\begin{Verbatim}[fontsize=\small,samepage=true,frame=none]
>>> print(df.lightsaber.head())
0      red
1     blue
2      red
3    green

>>> lb.fit_transform(df.lightsaber)[0:4]
array([[0, 0, 1],
       [1, 0, 0],
       [0, 0, 1],
       [0, 1, 0]])
\end{Verbatim}
\normalsize

You'll have to stick all your transformed variables back together with
\texttt{np.concatenate()} before passing them to a classifier.

\subsection{Comparing kNN with Naive Bayes}

Let's use the scikit-learn library to compare 3-NN with Naive Bayes. First, we
load the data and split it randomly into training and test sets:

\small
\begin{Verbatim}[fontsize=\small,samepage=true,frame=none]
w = pd.read_csv("warriors.csv")
train = w.sample(frac=.7)
test = w.drop(train.index)
\end{Verbatim}
\normalsize

Next, we one-hot encode the features (all of which are categorical):

\small
\begin{Verbatim}[fontsize=\small,samepage=true,frame=none]
def one_hot_this_puppy(df):
    lb = LabelBinarizer()
    lightsaber = lb.fit_transform(df.lightsaber)
    height = lb.fit_transform(df.height)
    beard = lb.fit_transform(df.beard)
    demeanor = lb.fit_transform(df.demeanor)
    return np.concatenate([lightsaber, height, beard, demeanor], axis=1)

train_features = one_hot_this_puppy(train)
test_features = one_hot_this_puppy(test)
\end{Verbatim}
\normalsize

We'll call scikit-learn's Naive Bayes implementation:

\small
\begin{Verbatim}[fontsize=\small,samepage=true,frame=none]
from sklearn.naive_bayes import MultinomialNB
naive_bayes = MultinomialNB()
naive_bayes.fit(train_features,train.type.as_matrix())
nb_predictions = naive_bayes.predict(test_features)
\end{Verbatim}
\normalsize

and its KNN implementation:

\small
\begin{Verbatim}[fontsize=\small,samepage=true,frame=none]
from sklearn.neighbors import KNeighborsClassifier
knn = KNeighborsClassifier(n_neighbors=3)
knn.fit(train_features, train.type.as_matrix())
knn_predictions = knn.predict(test_features)
\end{Verbatim}
\normalsize

Finally, print the results:

\small
\begin{Verbatim}[fontsize=\small,samepage=true,frame=none]
>>> print("Naive Bayes got {:.2f}% correct.".format(
    sum(nb_predictions == test.type)/len(test)*100))
>>> print("kNN(3) got {:.2f}% correct.".format(
    sum(knn_predictions == test.type)/len(test)*100))

Naive Bayes got 79.00% correct.
kNN(3) got 86.33% correct.
\end{Verbatim}
\normalsize

Pretty close, but for this data set, 3-NN has the edge.
\pagebreak

It's also illustrative to examine \textit{which} points Naive Bayes called
correctly but 3-NN missed, and vice versa.

\footnotesize
\begin{Verbatim}[fontsize=\small,samepage=true,frame=none]
nb_right = []
knn_right = []
for i in range(len(nb_predictions)):
    if nb_predictions[i] == test.type.iloc[i] != knn_predictions[i]:
        nb_right.append(i)
    if nb_predictions[i] != test.type.iloc[i] == knn_predictions[i]:
        knn_right.append(i)

nb_right_examples = test.iloc[nb_right]
knn_right_examples = test.iloc[knn_right]

>>> print("Naive Bayes got these right:")
>>> print(nb_right_examples)
>>> print("\nkNN got these right:")
>>> print(knn_right_examples)
\end{Verbatim}

\scriptsize
\begin{multicols}{2}
\begin{Verbatim}[fontsize=\small,samepage=true,frame=none]
Naive Bayes got these right:
    beard  height  demeanor lightsaber  type
7      no  medium      calm      green  Jedi
63     no   short  menacing      green  Jedi
116    no  medium      calm      green  Jedi
154    no  medium      calm      green  Jedi
185    no   short  menacing      green  Jedi
190    no   short  menacing      green  Jedi
279    no  medium      calm       blue  Jedi
600    no  medium      calm       blue  Jedi
827    no   short  menacing      green  Jedi
842    no   short  menacing      green  Jedi
909    no   short  menacing      green  Jedi
\end{Verbatim}
\columnbreak
\begin{Verbatim}[fontsize=\small,samepage=true,frame=none]
kNN got these right:
    beard  height  demeanor lightsaber  type
1      no   short      calm       blue  Sith
33     no   short      calm       blue  Sith
55     no   short      calm       blue  Sith
94     no  medium      calm      green  Sith
106    no   short      calm       blue  Sith
254    no   short  menacing      green  Sith
283    no   short      calm       blue  Sith
366    no  medium      calm       blue  Sith
380    no  medium      calm      green  Sith
385    no  medium      calm       blue  Sith
431    no   short      calm       blue  Sith
468    no   short      calm       blue  Sith
570    no   short      calm       blue  Sith
581    no   short      calm       blue  Sith
599    no   short      calm      green  Sith
686    no  medium      calm      green  Sith
713    no  medium      calm       blue  Sith
733    no   short      calm       blue  Sith
743    no  medium      calm       blue  Sith
815    no   short  menacing      green  Sith
870    no   short      calm       blue  Sith
931    no  medium      calm       blue  Sith
966    no  medium      calm      green  Sith
991    no  medium      calm       blue  Sith
\end{Verbatim}
\end{multicols}
\normalsize

Interestingly, all of Naive Bayes's errors were in one direction (incorrect
\texttt{Jedi} call) and all 3-NN's in the other. Naive Bayes seems to have the
most trouble with clean-shaven, calm, non-red-lightsabered Sith Lords. It
tends to call them ``Jedi," because two of the features --
\texttt{demeanor=calm} and \texttt{lightsaber=blue/green} -- are highly
\texttt{Jedi}-indicative. When factored in with just one
\texttt{Sith}-indicative feature (\texttt{beard=no}), those two features
overcome the third and tilt the Naive Bayes scales (incorrectly) towards Jedi.
3NN is able to pick out these \texttt{Sith} from the \texttt{Jedi} masses.

\subsection{Varying $k$}

One question that arises is: ``what's the
optimal value of $k$?" It may or may not surprise you to learn that there's no
easy answer to this, and that it varies widely by data set and in fact even by
the training set chosen.

Here's some code to try all (odd) values of $k$ between 1 and 301:
\small
\begin{Verbatim}[fontsize=\small,samepage=true,frame=none]
corrects_train = []
corrects_test = []
for k in np.arange(1,301,2):
    knn = KNeighborsClassifier(n_neighbors=k)
    knn.fit(train_features, train.type.as_matrix())
    knn_predictions_train = knn.predict(train_features)
    knn_predictions_test = knn.predict(test_features)
    corrects_train.append(sum(knn_predictions_train == train.type)/len(train)*100)
    corrects_test.append(sum(knn_predictions_test == test.type)/len(test)*100)
\end{Verbatim}
\normalsize

The plot looks like this:

\begin{center}
\includegraphics[width=0.9\textwidth]{k_vs_performance.pdf}
\end{center}

The performance on the training data is \textit{usually} better, as expected,
although you'll see some values of $k$ here in the 80-100 range where
paradoxically that's not true. More importantly, you'll notice that the line
shows no simple pattern. Very low numbers are good, and then performance tails
off...but then it goes back up again when $k$ is around 100 or so. Moreover,
\textit{the apparent optimal choice of k depends on the training set.} Here
are some more results from the same code, with different randomly-chosen
training sets:

\begin{center}
\includegraphics[width=0.31\textwidth]{k_vs_performance1.pdf}
\includegraphics[width=0.31\textwidth]{k_vs_performance2.pdf}
\includegraphics[width=0.31\textwidth]{k_vs_performance3.pdf}
\end{center}

We see some similarities here, but by no means perfect uniformity. This is
because a kNN classifier is a high-variance classifier (especially for low
$k$), as described in the next section.

\section{The ``Bias-Variance Tradeoff"}

One of the most deep and fundamental principles of machine learning goes by
the name of \textbf{the bias-variance tradeoff}. It's a version of the ``no
free lunch" principle: classifiers that possess one advantage inherently have
a counteracting disadvantage.

Here are the two key terms:

\begin{itemize}
\itemsep.1em
\item \textbf{bias}\footnote{Tip: when thinking about this topic, try to erase the previous definitions you
have for the word ``\textit{bias}" that you have running around in your brain.
They probably don't apply. When I think of ``bias" I think of a systematic
propensity to err \textit{in a particular direction}. Like, ``you're biased
towards thinking people are shorter than they are, because you're tall." But
\textit{bias} here just means ``being `off' in general, in whatever direction,
and probably not the same direction each time."}
 -- the degree to which a model is expected to misclassify
items, even those in its own training set, because it just can't ``bend" 
enough to catch all the detail.
\item \textbf{variance} -- the degree to which a model will differ based on
the particular training set that is chosen, because it ``bends" too much to
try and catch every data point.
\end{itemize}

The basic theory is as follows. Suppose you have a very simple classifier that
tries to identify a single, smooth, perfectly straight line separating the
data points. Maybe our model says ``just use the \texttt{lightsaber} feature
by itself, and call `\texttt{Jedi}' or `\texttt{Sith}' based on whether the
test point's lightsaber value is more typical of the \texttt{Jedi} Knights and
\texttt{Sith} Lords in the data set.

This simple classifier can be said to have \textbf{low variance}. Why? Because
if, instead of those particular 714 training points chosen from the
population, I had sampled a different 714 training points, \textit{the
classification decisions I would make on new data would be about the same.} To
see why this is the case, consider how different the training sets would have
to me in order for me to make any \textit{different} decisions. Let's say my
original 714 training points had:

\begin{center}
\begin{Verbatim}[fontsize=\small,samepage=true,frame=none]
        Jedi    Sith
red       17      85
blue     283      28
green    299       2
\end{Verbatim}
\end{center}

for their lightsaber colors. Clearly, my very simply classifier is going to
predict \texttt{Sith} for any \texttt{red}-lightsaber wielder and
\texttt{Jedi} for anyone else. Now if I had collected a different training
set, the numbers above would be different, but they would have to be a
\textit{lot} different to change any of my future predictions. I'd have to
have an additional 68 \texttt{red} \texttt{Jedi}, for instance, or 255
additional \texttt{blue} \texttt{Sith}, or a whopping 297 additional
\texttt{green} \texttt{Sith}, without any increase in the other class. Not
likely.

That's the bright side. On the dark side, by classifying points in this way,
this simple model is going to get a lot of data points wrong: $17+28+2=47$
points incorrect \textit{in its own training set}, in fact. This mean it's a
pretty \textbf{high bias} classifier.

It's a tradeoff. In exchange for being virtually immune to the idiosyncrasies
in the training sample, we've sacrificed some precision.

The other extreme can be seen by visualizing a 1-NN classifier. Now, for a
future data point, we're going to find the single nearest neighbor to it in
our training set, and call \texttt{Jedi} or \texttt{Sith} accordingly. This is
very \textbf{low bias} -- we're guaranteed to get every training point
correct, in fact!\footnote{Unless there are data points with identical feature
values but different labels, in which case whatcha gonna do.} But it is also
very \textbf{high variance}: a different training set will result in many,
many different classifications. Our classifier has essentially memorized a
zillion individual examples, and will regurgitate all that intricate detail on
command. But had we gotten a different random sample, the ``single nearest
neighbor" to a test point would have been totally different, and might result
in a different classification.
     
There are formulas you can use to work out quantitatively what the estimated
bias and variance are for a classifier. But I think this dichotomy is more
important from a conceptual perspective. The main lesson is: the more your
classifier ``bends" (\textit{i.e.}, has a very wiggly and unpredictable
boundary between its ranges of predicted classes) the more precise it
potentially is in selectively labeling data points. However, this also means
that a lot of its decisions are going to hinge on the particular data sample
it was trained on, which means it will probably be picking up noise in
addition to signal.


\section{The ``Curse of Dimensionality"}

Another key ML principle goes by the humorous name \textbf{the curse of
dimensionality}. It goes something like this. Suppose that in addition to
\texttt{height}, \texttt{demeanor}, \texttt{lightsaber}, and \texttt{beard},
we also have at our disposal many other features about Jedi \& Sith: their
ages, salaries, home planets, personality profiles, number of traffic tickets,
\textit{etc.} This is good news, right? It gives us more and more data to sift
through statistically and discover correlations, right? What's not to like?

Yes, this is potentially good news, but with a catch. It turns out that as we
increase the number of features more and more (or, put another way, as we
increase the ``dimensionality" of the problem from our original four
dimensions to 5, 6, 7, ..., and up into the hundreds) we have a data sparsity
problem. In terms of kNN classifiers, you can think of it this way:
effectively, \textit{all} the points effectively become far apart. We want to
find the 3 (or 15, or 77) ``nearest neighbors," but when we have so many
features (many of which will turn out to be irrelevant for the problem), we
find that \textit{any} two data points will differ on a large number of them.
Essentially, our noise swamps our signal.

Another aspect of the problem concerns the amount of training data that we
need. When we have tons of features, we need an enormous number of training
samples in order to guarantee that we have several samples with each
combination of values (or their ranges). In our original problem, we didn't
need to much data to ensure that we'd have at least a few \texttt{tall},
\texttt{green}-lightsabered, \texttt{menacing} individuals with \texttt{no}
beard. But how many training points will we need before we have several
\texttt{tall}, \texttt{green}-lightsabered, \texttt{menacing}, 38-year-old,
extroverted, abstract-thinking individuals from Tatooine with 4 traffic
tickets and \texttt{no} beard who earn 6,000 Galactic credits annually?
Answer: a lot. And if we don't have that many, a kNN classifier trained on all
those features is going to suffer trying to find ``people like this one" in
order to make a prediction.

The moral to the story is that we will often need to trim our set of features,
sometimes aggressively, if we have too many. It feels wasteful -- all this
great data on all these people, and we're throwing it away! -- but it's the
right thing to do. This process is called \textbf{feature selection}, and
we'll cover it more fully in a future lesson.

\vspace{.1in}
In the meantime, here are three useful principles for avoiding the curse:

\begin{itemize}
\itemsep.1em
\item By far the most important thing to recognize is that if you have
additional features that are \textit{truly} (and pretty strongly) correlated
with the target, it's okay to include them. Adding more stuff that actually
contains useful information for prediction isn't generally a problem. The
problem is when you add features that are largely \textit{noise}, since
they will overwhelm the truly significant features if there are too many of
them.
\item A common rule of thumb is to make sure you have at least \textit{five}
training points for each combination of feature values. (In our Jedi example,
with three \texttt{height}s and three \texttt{lightsaber} values and two
values for the other features, this would imply a training set of at least $(3
\times 3 \times 2 \times 2) \times 5 = 180$ training points.)
\item When the size of the training set becomes too low, use
\textit{parametric} methods (like Naive Bayes or linear regression) rather
than instance-based methods (like kNN) because they suffer less from the
curse.
\end{itemize}


